\documentclass[10pt,a4paper]{amsart}
\usepackage[margin=19mm]{geometry}
\usepackage{amsmath,amssymb,mathtools}
\usepackage[scr=rsfs,cal=euler]{mathalfa}
\usepackage{xfrac}
\usepackage[unicode,hidelinks,bookmarks=false]{hyperref}
\newcommand{\ie}{i.e.\ }
\newcommand{\cf}{cf.\ }
\newcommand{\resp}{respectively\ }
\newcommand{\Subsection}{\subsection}
\providecommand{\nobreakdash}{\nobreak\hbox{-}}
\setlength{\emergencystretch}{2em}
\multlinegap0pt
\usepackage{amssymb}
\newtheorem{definition}{Definition}
\newtheorem{theorem}{Theorem}
\title{Criteria for existence of semigroup homomorphisms and projective rank functions}
\author{George M. Bergman}
\date{Source: arXiv:2603.20628v1 (CC BY 4.0)}
\begin{document}
\maketitle

\noindent\emph{Source and licence.} Criteria for existence of semigroup homomorphisms and projective rank functions. Version arXiv:2603.20628v1 (CC BY 4.0), distributed by arXiv under the Creative Commons Attribution 4.0 International licence, http://creativecommons.org/licenses/by/4.0/, as recorded in the arXiv licence metadata for this exact version and shown on the abstract page https://arxiv.org/abs/2603.20628v1. Full licence notice: \texttt{licenses/arxiv-2603.20628-CC-BY-4.0.txt}.\par\smallskip

\noindent\textit{Editorial scope.} Counted unit: the article's theorem T.main (source0.tex lines 140--180). The article's complete original proof of it is delivered with it (source0.tex lines 182--421, one \texttt{proof} environment of 240 lines). No mathematics is written, repaired or re-derived: the statement and the proof are the article's own lines, in the article's own notation, and no retained line is substituted or re-flowed.  Retained uncounted as the source-local context the proof needs: lines 109--122.  Nothing was omitted from any retained span, so the package carries no omission record.  No retained line contains a citation, so the wrapper carries no bibliography and no bibliography entry.  No retained line contains a figure environment, an image-inclusion command or drawing code, so no image is delivered and the package declares no asset dependency.  Editorial formatting only, in addition: an AMS \texttt{amsart} preamble that restates the article's own declarations and macros used by the retained lines, namely the environments \texttt{definition}, \texttt{theorem} on separate counters, together with the standard packages those lines need.  The retained environments print in this document as Definition 1, Theorem 1 in order of appearance, whereas the article prints the same items with the numbering of its own full document.  The complete native source of the article as distributed in the arXiv e-print for this exact version is bundled unchanged as \texttt{sources/196864/source0.tex}.
\begin{definition}\label{D.divisor}
If $S$ is a semigroup, $S^1$ will denote the
monoid obtained by adjoining an identity element to $S.$

An element $s$ of a semigroup $S$ will be called
a {\em divisor} of an element $t\in S$
if $t = a\,s\,b$ for some $a, b\in S^1.$

An element $s$ of a semigroup $S$ will be called
a {\em weak divisor} of a
subset $A\subseteq S$ if there exists a positive
integer $d$ such that $s^d$ is a divisor of some element
of $A^d$ {\rm(}i.e., of some product $a_1\dots\,a_d$ with $a_i\in A).$
\end{definition}

\begin{theorem}\label{T.main}
Let $P,$ $S$ and $T$ be semigroups, $f:P\to S$ and $g:P\to T$
semigroup homomorphisms, and $X$ a generating set for $S.$
Suppose that
%
\begin{equation}\begin{minipage}[c]{35pc}\label{d.oAdiv}
every element of $S$ is a divisor of some element of $f(P),$
\end{minipage}\end{equation}
%
\begin{equation}\begin{minipage}[c]{35pc}\label{d.cancel}
$T$ is right and left cancellative {\rm(}i.e., for $x,y\in T$
and $a,b\in T^1,$ $a\,x\,b\ = a\,y\,b \implies x = y),$
\end{minipage}\end{equation}
%
\begin{equation}\begin{minipage}[c]{35pc}\label{d.pwr}
$T$ is power cancellative
{\rm(}i.e., $x^d = y^d \implies x = y$ for $x,\,y\in T,$ $d>0),$
\end{minipage}\end{equation}
%
and
%
\begin{equation}\begin{minipage}[c]{35pc}\label{d.fin_wk_div}
every finite subset of $g(P) \subseteq T$ has
only finitely many weak divisors in $T$
{\rm(}Definition~\ref{D.divisor} above{\rm)}.
\end{minipage}\end{equation}

Then the following conditions are equivalent:
%
\begin{equation}\begin{minipage}[c]{35pc}\label{d.realization}
For every relation $f(p) = w(x_1,\,\dots\,,\,x_n)$ holding
in $S,$ with $p\in P,$ $w$ a semigroup word, and
$x_1,\,\dots\,,\,x_n\in X,$ there exist
$t_1,\,\dots\,,\,t_n\in T$ satisfying $g(p) = w(t_1,\,\dots\,,\,t_n).$
\end{minipage}\end{equation}
%
\begin{equation}\begin{minipage}[c]{35pc}\label{d.oEh}
There exists a homomorphism $h:S\to T$ such that $g=hf.$
\end{minipage}\end{equation}
%
\end{theorem}

\begin{proof}
Clearly,~\eqref{d.oEh} implies~\eqref{d.realization}
(with $t_i=h(x_i)).$
So what we must prove is the reverse implication.

Observe that a homomorphism $h:S\to T$ is determined by
its restriction to $X;$ and clearly, given any set map
%
\begin{equation}\begin{minipage}[c]{35pc}\label{d.eta}
$\eta:\,X\to T,$
\end{minipage}\end{equation}
%
the necessary and sufficient condition for $\eta$ to
extend to a homomorphism $S\to T$ is that
%
\begin{equation}\begin{minipage}[c]{35pc}\label{d.w12}
for every relation
$w_1(x_1,\,\dots\,,\,x_m) = w_2(x_1,\,\dots,\,x_m)$
holding in $S,$ where $w_1$ and $w_2$ are semigroup words
and $x_1,\,\dots\,,\,x_m\in X,$ the corresponding relation
$w_1(\eta(x_1),\,\dots\,,\,\eta(x_m)) =
w_2(\eta(x_1),\,\dots,\,\eta(x_m))$ holds in~$T.$
\end{minipage}\end{equation}

Given homomorphisms $f:P\to S$ and $g:P\to T$ as in
the statement of the theorem, and a map~\eqref{d.eta},
when will that map determine a homomorphism $h$ such that $g=hf$?
I claim that if
%
\begin{equation}\begin{minipage}[c]{35pc}\label{d.w=f}
for every relation $f(p)=w(x_1,\,\dots\,,\,x_n)$
holding in $S,$ with $p\in P$ and $x_1,\,\dots,\,x_n\in X,$ one has
$g(p)=w(\eta(x_1),\,\dots,\,\eta(x_n))$ in $T,$
\end{minipage}\end{equation}
%
then assuming~\eqref{d.oAdiv}-\eqref{d.fin_wk_div},
this will imply that~\eqref{d.w12} holds,
yielding a homomorphism $h:S\to T;$ and clearly,~\eqref{d.w=f}
then shows that $h$ will satisfy the desired relation $g=hf.$

To prove~\eqref{d.w12} assuming~\eqref{d.w=f}, consider
any relation as in the hypothesis of~\eqref{d.w12}.
Let us write $s$ for the common value in $S$
of the two sides of this relation.
Then by~\eqref{d.oAdiv},
$s$ is a divisor of $f(p)$ for some $p\in P.$
Expressing the right and left factors that
carry $s$ to $f(p)$ in terms of
the generating set $X$ (using a list $x_1,\,\dots\,,\,x_n\in X$
extending $x_1,\,\dots\,,\,x_m,$ and allowing empty words
if either or both factors are $1),$ we get a relation
%
\begin{equation}\begin{minipage}[c]{35pc}\label{d.uw12v}
$\!\!u(x_1,\dots,\,x_n)\,w_1(x_1,\dots,\,x_m)\,
v(x_1,\dots,\,x_n) = f(p) =
u(x_1,\dots,\,x_n)\,w_2(x_1,\dots,\,x_m)\,v(x_1,\dots,\,x_n)\!\!$
\end{minipage}\end{equation}
%
in $S.$
Regarding this as two equations of the sort appearing
in the hypothesis of~\eqref{d.w=f}, applying~\eqref{d.w=f}
to each, and, since the resulting left-hand sides
are both $g(p),$ equating the right-hand sides,
we get an equation in $T$ which,
by the cancellativity condition~\eqref{d.cancel},
yields the conclusion of~\eqref{d.w12}, as desired.

So,
%
\begin{equation}\begin{minipage}[c]{35pc}\label{d.iff}
To prove the Theorem, it will suffice to show that
assuming~\eqref{d.realization},
there must exist a map~\eqref{d.eta} for which~\eqref{d.w=f} holds.
\end{minipage}\end{equation}
%

In constructing such an $\eta,$
the following notation will be useful:
%
\begin{equation}\begin{minipage}[c]{35pc}\label{d.Sigma}
We shall denote by $\Sigma$ the family of equations
$g(p)=w(\eta(x_1),\,\dots,\,\eta(x_n))$ arising
(as in the conclusion of~\eqref{d.w=f})
from all relations $f(p)=w(x_1,\,\dots\,,\,x_n)$ that
hold in $S,$ with the symbols $\eta(x_1),\,\dots,\,\eta(x_n)$
now regarded as an $\!n\!$-tuple of $\!T\!$-valued unknowns
(though the symbols $g(p)$ will, as in~\eqref{d.w=f},
still denote the indicated constants in~$T).$
\end{minipage}\end{equation}

Below, we shall first deduce from~\eqref{d.realization} that
every {\em finite} subset of $\Sigma$ has a solution
in $T^X,$ then show from this that the whole family $\Sigma$
has such a solution.

To get the result on finite subsets of $\Sigma,$
let us start by showing that given any two equations in $\Sigma,$
%
\begin{equation}\begin{minipage}[c]{35pc}\label{d.uu'}
$g(p)=u$ \ and \ $g(p')=u',$
\end{minipage}\end{equation}
%
there exists a nonnegative integer $d$ such that the equation
%
\begin{equation}\begin{minipage}[c]{35pc}\label{d.u^du'}
$g(p^d\,p')\ =\ u^d\,u',$
\end{minipage}\end{equation}
%
which also belongs to $\Sigma,$ has the same solution-set
in $T^X$ as the pair of equations~\eqref{d.uu'}.

For this purpose, let $\eta(x_1),\dots,\eta(x_n)$ be the finitely many
variables which appear in one or both of the words
$u$ and $u',$ and consider
assignments of values to $\eta(x_1),\dots,\eta(x_n)$ such
that~\eqref{d.u^du'} holds for at least one value of~$d.$
Every value of $u(\eta(x_1),\dots,\eta(x_n))$
that such an assignment leads to will, in view of~\eqref{d.u^du'},
be a weak divisor (Definition~\ref{D.divisor})
of the $\!2\!$-element set $\{g(p),\ g(p)\,g(p')\},$
so by~\eqref{d.fin_wk_div},
%
\begin{equation}\begin{minipage}[c]{35pc}\label{d.fin_u-vals}
there are only finitely many values in $T$ that
$u(\eta(x_1),\dots,\eta(x_n))$ can have in solutions
to~\eqref{d.u^du'}, as $d$ ranges over all nonnegative integers.
\end{minipage}\end{equation}

Moreover, I claim that for any such value of
$u(\eta(x_1),\dots,\eta(x_n))$ that is
{\em not} equal to~$g(p),$ the equality~\eqref{d.u^du'} can only
hold for one value of~$d.$
Indeed, assuming that for some $d$ and some
choice of values for $\eta(x_1),\dots,\eta(x_n)$
the equation~\eqref{d.u^du'}
holds, with $g(p)\neq u,$ we see by power-cancellativity
of $T$~\eqref{d.pwr} that %
that choice will make $g(p)^c\neq u^c$ for every $c>0,$ so
if we multiply this inequality on the right by~\eqref{d.u^du'},
right cancellativity of $T$ (condition~\eqref{d.cancel}) shows
that $g(p)^{c+d}g(p')\neq u^{c+d}\,u'.$
So taking for $d$ the least value for which the given
choice of values for $\eta(x_1),\dots,\eta(x_n)$
make~\eqref{d.u^du'} hold, we see that it is the only such value.
Hence, as $u(\eta(x_1),\dots,\eta(x_n))$ ranges over the
finitely many values referred to
in~\eqref{d.fin_u-vals}, only finitely many values of $d$
allow solutions of~\eqref{d.u^du'} with
$g(p)\neq u(\eta(x_1),\dots,\eta(x_n)).$

So let us choose a $d$ that is {\em not} one of those finitely
many values; so that all solutions of~\eqref{d.u^du'}
for that choice of $d$ satisfy $g(p)=u.$
Then by left cancellativity (hypothesis~\eqref{d.cancel}),
those solutions also satisfy $g(p')=u'.$
Thus, for that choice of $d,$ the values of
$\eta(x_1),\dots,\eta(x_n)$ that make~\eqref{d.u^du'} hold are, as
claimed, precisely those that make both equalities of~\eqref{d.uu'}
hold.

Repeatedly applying this result, we see that the solution-set
in $T^X$ of any {\em finite} subset of $\Sigma$ is
equal to the solution-set of a single element of $\Sigma.$
By~\eqref{d.realization}, every member of $\Sigma$
has a nonempty solution-set; hence the above argument
shows that every finite subset of $\Sigma$ does.

Finally, let us show from this that the whole set $\Sigma$
has a nonempty solution-set.

We begin by noting that by~\eqref{d.oAdiv}, each $x\in X$ is
a divisor of some element $f(p_x)$ $(p_x\in P),$ hence
%
\begin{equation}\begin{minipage}[c]{35pc}\label{d.u_x}
for each $x\in X$
there is some member of $\Sigma,$ say $g(p_x) = u_x,$ such that
the word $u_x$ involves the variable $\eta(x).$
\end{minipage}\end{equation}

For each $x\in X,$ let us fix such an equation $g(p_x)=u_x$ in $\Sigma,$
and denote by $T_x$ the set of divisors in $T$ of the element $g(p_x).$
By~\eqref{d.fin_wk_div} each $T_x$ is finite
(since every divisor of an element $g(p_x)$ is, in particular,
a weak divisor of the singleton set $\{g(p_x)\}).$
Thus, in any solution in the semigroup $T$ of any subsystem
of $\Sigma$ which includes the equation $g(p_x) = u_x,$ the
value given to $\eta(x)$ must be a member of the finite set~$T_x.$

Now let $\Sigma_0$ be any finite subset of $\Sigma,$ let
$\eta(x_1),\dots,\eta(x_n)$ be the variables occurring in
the right-hand sides of the equations in $\Sigma_0,$ and
let $\Sigma_1$ be the finite set obtained by adjoining to
$\Sigma_0$ the additional
equations $g(p_{x_i})=u_{x_i}$ for $i=1,\dots,n.$
Since $\Sigma_1$ is a finite subset
of $\Sigma,$ it will, as we have shown, have
solutions; and since it contains the equations
$g(p_{x_i}) = u_{x_i}$ for $i=1,\dots,n,$
those solutions will have $\!x_i\!$-coordinate in $T_{x_i}$
for $i=1,\dots,n.$
But solutions to $\Sigma_1$ are
also solutions to its subfamily $\Sigma_0,$ so
there are solutions to the latter family
with $\!x_i\!$-coordinate in $T_{x_i}$ for $i=1,\dots,n.$
And since the equations in $\Sigma_0$ involve no variables
but $\eta(x_1),\dots,\eta(x_n),$ we can, in such solutions,
modify the values assigned to all other coordinates
$\eta(x)$ $(x\in X)$
in any way; in particular, replace each such coordinate $\eta(x)$
by a member of the corresponding set $T_x.$
This shows that
%
\begin{equation}\begin{minipage}[c]{35pc}\label{d.in_prod_T_i}
the solution-set in $T_x$
of every finite subset $\Sigma_0\subseteq\Sigma$
has nonempty intersection with $\prod_{x\in X} T_x.$
\end{minipage}\end{equation}

Let us now regard $\prod_X\,T_x$
as a compact topological space, under the product of the
discrete topologies on the finite factors $T_x,$ and consider each
semigroup word in the variables $\{\eta(x)\,|\,x\in X\}$ as defining a
map $\prod_X\,T_x\to T.$
Regarding $T$ as a discrete space,
each such map is continuous, since it depends on only
finitely many coordinates of $\prod_X\,T_x,$ so the
set of solutions in $\prod_X\,T_x$ to
each equation in $\Sigma$ is closed.

The class of finite subsets of $\Sigma$ is closed
under finite unions, hence the class of their solution-sets in
$\prod_X\,T_x$ is closed under finite intersections.
By the preceding results, all such finite intersections are
nonempty closed sets,
hence the compactness of $\prod_X\,T_x$ shows that the solution-set of
the full family $\Sigma$ is nonempty.
A member of this intersection will be a map~\eqref{d.eta}
satisfying~\eqref{d.w=f},
which by~\eqref{d.iff} completes the proof of the Theorem.
\end{proof}

\end{document}
